\documentclass[12pt,a4paper]{article}

% ─── Packages ───────────────────────────────────────────────────────────────
\usepackage[utf8]{inputenc}
\usepackage[T1]{fontenc}
\usepackage{amsmath, amssymb, amsthm}
\usepackage{physics}          % \ket, \bra, \expval, \comm
\usepackage{braket}
\usepackage{graphicx}
\usepackage{hyperref}
\usepackage{cleveref}
\usepackage{geometry}
\usepackage{xcolor}
\usepackage{booktabs}
\usepackage{microtype}
\usepackage[numbers, sort&compress]{natbib}

\geometry{margin=2.5cm}

\hypersetup{
  colorlinks = true,
  linkcolor  = blue!70!black,
  citecolor  = blue!70!black,
  urlcolor   = blue!70!black
}

% ─── Custom commands ─────────────────────────────────────────────────────────
\newcommand{\Hil}{\mathcal{H}}
\newcommand{\Lind}{\mathcal{L}}
\newcommand{\Tr}{\operatorname{Tr}}
\newcommand{\ii}{\mathrm{i}}
\newcommand{\ee}{\mathrm{e}}
\newcommand{\wq}{\omega_q}
\newcommand{\wt}{\omega_{\mathrm{TLS}}}
\newcommand{\wd}{\omega_d}
\newcommand{\D}{\Delta}
\newcommand{\smq}{\sigma_-^{(q)}}
\newcommand{\spq}{\sigma_+^{(q)}}
\newcommand{\smt}{\sigma_-^{(t)}}
\newcommand{\spt}{\sigma_+^{(t)}}
\newcommand{\szq}{\sigma_z^{(q)}}
\newcommand{\szt}{\sigma_z^{(t)}}

% ─── Theorem environments ────────────────────────────────────────────────────
\newtheorem{remark}{Remark}

% ════════════════════════════════════════════════════════════════════════════
\begin{document}

\title{%
  \textbf{Full Quantum Evolution of a Qubit Coupled to a Two-Level System}\\[0.4em]
  \large From the Jaynes--Cummings Model to the Lindblad Master Equation\\[0.4em]
  \normalsize Theoretical Background for TLS-QEC Thesis
}
\author{Ege Eroglu\\[0.2em]
  \small Supervisors: Mathieu Fechant, Nico Gosling\\[0.2em]
  \small Physikalisches Institut, KIT Karlsruhe
}
\date{\today}
\maketitle
\tableofcontents
\newpage

% ════════════════════════════════════════════════════════════════════════════
\section{Introduction and Physical Motivation}

Superconducting qubits inevitably interact with material defects in their
substrate and tunnel junctions. These defects are well modelled as
\emph{two-level systems} (TLS): quantum objects with a ground state
$\ket{g}$ and an excited state $\ket{e}$, just like the qubit itself
\cite{Muller2019,Lisenfeld2019}.

The interaction between a qubit and a TLS introduces several detrimental
effects. First, it modifies the qubit transition frequency through virtual
photon exchange. Second, it opens an additional decay channel: the qubit
can relax by transferring energy to the TLS, which subsequently releases it
back, leading to \emph{non-exponential} and \emph{non-Markovian} decay
\cite{Ashhab2006}. Third, and most relevant to quantum error correction
(QEC), the qubit--TLS interaction introduces \emph{time-correlated errors}
that are known to degrade surface code performance \cite{Kam2025}.

The goal of this document is to derive, step by step, the full quantum
model of the qubit--TLS system, starting from first principles and arriving
at the Lindblad master equation whose solutions will be compared with the
Solomon equations in subsequent work.

% ════════════════════════════════════════════════════════════════════════════
\section{The Hilbert Space and Bare Hamiltonians}

\subsection{Hilbert space structure}

The total Hilbert space of the qubit--TLS system is the tensor product
\begin{equation}
  \Hil = \Hil_q \otimes \Hil_{\mathrm{TLS}},
  \qquad
  \dim \Hil = 4.
\end{equation}
Each subsystem is a two-dimensional space spanned by
$\{\ket{g},\ket{e}\}$. The four computational basis states of the
composite system are
\begin{equation}
  \ket{ee},\quad \ket{eg},\quad \ket{ge},\quad \ket{gg},
\end{equation}
where the first (second) label refers to the qubit (TLS).

\subsection{Individual Hamiltonians}

The free Hamiltonian of each two-level system is written in terms of the
Pauli $z$-operator ($\hbar = 1$ throughout)
\begin{align}
  H_q    &= \frac{\wq}{2}\,\szq, \label{eq:Hq}\\[4pt]
  H_{\mathrm{TLS}} &= \frac{\wt}{2}\,\szt, \label{eq:Ht}
\end{align}
where $\wq$ ($\wt$) is the transition frequency of the qubit (TLS), and
$\sigma_z = \ket{e}\bra{e} - \ket{g}\bra{g}$.
The raising and lowering operators are
$\sigma_+ = \ket{e}\bra{g}$, $\sigma_- = \ket{g}\bra{e}$,
satisfying $[\sigma_+,\sigma_-] = \sigma_z$.

% ════════════════════════════════════════════════════════════════════════════
\section{The Interaction Hamiltonian}

\subsection{Full dipole coupling}

The most general bilinear coupling between the two spins, arising from
capacitive or inductive coupling in circuit QED \cite{Blais2021}, is
\begin{equation}
  H_{\mathrm{int}}^{\mathrm{full}}
  = g\,\sigma_x^{(q)} \otimes \sigma_x^{(t)}
  = g\,(\spq + \smq)(\spt + \smt).
\end{equation}
Expanding, this contains four terms:
\begin{equation}
  H_{\mathrm{int}}^{\mathrm{full}}
  = g\Bigl(
      \underbrace{\spq\smt + \smq\spt}_{\text{energy-conserving}}
    + \underbrace{\spq\spt + \smq\smt}_{\text{counter-rotating}}
    \Bigr).
  \label{eq:Hint_full}
\end{equation}

\subsection{Rotating-wave approximation}

In the interaction picture with respect to $H_0 = H_q + H_{\mathrm{TLS}}$,
the counter-rotating terms oscillate at frequency $\wq + \wt$, while the
energy-conserving terms oscillate at $|\wq - \wt| = |\D|$.
When
\begin{equation}
  g \ll \wq,\,\wt,
  \label{eq:RWA_condition}
\end{equation}
the counter-rotating terms average to zero over one oscillation period and
may be dropped. This is the \textbf{rotating-wave approximation} (RWA)
\cite{Haroche2006,WallsMilburn2008}. The resulting interaction Hamiltonian is
\begin{equation}
  \boxed{
  H_{\mathrm{int}} = g\bigl(\spq\smt + \smq\spt\bigr),
  }
  \label{eq:Hint_RWA}
\end{equation}
which conserves the total excitation number
$N_{\mathrm{exc}} = \frac{1}{2}(\szq + \szt + 2)$.

\begin{remark}
  In the \textbf{ultrastrong coupling} regime $g/\wq \gtrsim 0.1$, the RWA
  breaks down and the full Hamiltonian \eqref{eq:Hint_full} must be used.
  All results below assume $g \ll \wq$.
\end{remark}

% ════════════════════════════════════════════════════════════════════════════
\section{The Total Hamiltonian and Block Structure}

The total Hamiltonian in the lab frame under RWA is the
\textbf{Jaynes--Cummings-like} Hamiltonian \cite{Blais2021,Ashhab2006}
\begin{equation}
  \boxed{
  H = \frac{\wq}{2}\szq + \frac{\wt}{2}\szt
    + g\bigl(\spq\smt + \smq\spt\bigr).
  }
  \label{eq:H_total}
\end{equation}

\subsection{Block-diagonal structure}

Because $H$ conserves $N_{\mathrm{exc}}$, it is block diagonal in the
basis ordered by excitation number:
\begin{equation}
  H =
  \underbrace{
    \begin{pmatrix} \wq/2 + \wt/2 \end{pmatrix}
  }_{\substack{N=2\\\ket{ee}}}
  \oplus
  \underbrace{
    \begin{pmatrix}
      (\wq - \wt)/2 & g \\
      g & -(\wq-\wt)/2
    \end{pmatrix}
  }_{\substack{N=1\\\{\ket{eg},\ket{ge}\}}}
  \oplus
  \underbrace{
    \begin{pmatrix} -\wq/2 - \wt/2 \end{pmatrix}
  }_{\substack{N=0\\\ket{gg}}}.
  \label{eq:H_block}
\end{equation}
All non-trivial dynamics occur in the \textbf{single-excitation subspace}
$\mathcal{H}_1 = \mathrm{span}\{\ket{eg},\ket{ge}\}$.

% ════════════════════════════════════════════════════════════════════════════
\section{Diagonalization and Dressed States}

\subsection{Single-excitation subspace}

Defining the detuning $\D = \wq - \wt$, the Hamiltonian in
$\mathcal{H}_1$ reads
\begin{equation}
  H_1 = \frac{\D}{2}\sigma_z + g\,\sigma_x
      = \begin{pmatrix} \D/2 & g \\ g & -\D/2 \end{pmatrix}.
  \label{eq:H1}
\end{equation}

\subsection{Eigenvalues}

The eigenvalues of $H_1$ are
\begin{equation}
  E_\pm = \pm\frac{1}{2}\sqrt{\D^2 + 4g^2} \equiv \pm\Omega_R,
  \label{eq:eigenvalues}
\end{equation}
where $\Omega_R = \tfrac{1}{2}\sqrt{\D^2 + 4g^2}$ is the
\textbf{generalised Rabi frequency}.

\subsection{Eigenstates (dressed states)}

The eigenstates are
\begin{align}
  \ket{+} &= \cos\theta\,\ket{eg} + \sin\theta\,\ket{ge}, \label{eq:ep}\\
  \ket{-} &= -\sin\theta\,\ket{eg} + \cos\theta\,\ket{ge}, \label{eq:em}
\end{align}
with the \textbf{mixing angle}
\begin{equation}
  \tan 2\theta = \frac{2g}{\D},
  \qquad
  \cos 2\theta = \frac{\D}{\sqrt{\D^2+4g^2}},
  \qquad
  \sin 2\theta = \frac{2g}{\sqrt{\D^2+4g^2}}.
  \label{eq:mixing_angle}
\end{equation}

\subsection{Physical limits}

\paragraph{Resonance ($\D = 0$).}
The mixing angle is $\theta = \pi/4$ and the dressed states become
symmetric and antisymmetric superpositions
\begin{equation}
  \ket{\pm} = \frac{1}{\sqrt{2}}\bigl(\ket{eg} \pm \ket{ge}\bigr),
\end{equation}
with splitting $E_+ - E_- = 2g$ (the \textbf{vacuum Rabi splitting}).

\paragraph{Dispersive limit ($|\D| \gg g$).}
The eigenstates are nearly bare states with a small admixture $\approx g/\D$,
and the dressed state splitting is
\begin{equation}
  E_+ - E_- \approx |\D| + \frac{2g^2}{|\D|}.
\end{equation}
The second term is the \textbf{dispersive (Lamb) shift}, which displaces the
qubit frequency by $\chi = g^2/\D$ \cite{Blais2021}.

% ════════════════════════════════════════════════════════════════════════════
\section{Closed-System Time Evolution}

\subsection{Time-evolution operator}

For a time-independent Hamiltonian, the time-evolution operator is
\begin{equation}
  U(t) = \ee^{-\ii H t}.
\end{equation}
In the single-excitation subspace, this can be evaluated analytically
using the identity for a $2\times 2$ traceless Hamiltonian
$H_1 = \mathbf{h}\cdot\boldsymbol{\sigma}$ with
$|\mathbf{h}| = \Omega_R$:
\begin{equation}
  \ee^{-\ii H_1 t}
  = \cos(\Omega_R t)\,\mathbf{1}
  - \ii\sin(\Omega_R t)\,\frac{H_1}{\Omega_R}.
  \label{eq:U_analytic}
\end{equation}

\subsection{Rabi oscillations}

Starting in $\ket{\psi(0)} = \ket{eg}$ (qubit excited, TLS ground), the
probability of finding the qubit excited at time $t$ is
\begin{equation}
  \boxed{
  P_e(t) = 1 - \frac{g^2}{\Omega_R^2}\sin^2(\Omega_R t).
  }
  \label{eq:Rabi}
\end{equation}
At resonance ($\D=0$), this reduces to $P_e(t) = \cos^2(gt)$: the
excitation oscillates completely between qubit and TLS at the Rabi
frequency $g$ \cite{Haroche2006}.
Off resonance, the oscillations are faster ($\Omega_R > g$) but
incomplete (the qubit never fully de-excites).

\subsection{State vector at time $t$}

In the $\{\ket{eg},\ket{ge}\}$ subspace:
\begin{equation}
  \ket{\psi(t)} =
  \Bigl[\cos(\Omega_R t) - \ii\frac{\D}{2\Omega_R}\sin(\Omega_R t)\Bigr]\ket{eg}
  - \ii\frac{g}{\Omega_R}\sin(\Omega_R t)\,\ket{ge}.
  \label{eq:psi_t}
\end{equation}

% ════════════════════════════════════════════════════════════════════════════
\section{Open System: The Lindblad Master Equation}

A real qubit and TLS are embedded in a wider electromagnetic environment
and therefore decay. To account for dissipation, we promote the pure state
$\ket{\psi}$ to a \textbf{density matrix} $\rho \in \mathcal{B}(\Hil)$
and use the \textbf{Gorini--Kossakowski--Sudarshan--Lindblad} (GKSL) master
equation \cite{Breuer2002,Gardiner2004}:
\begin{equation}
  \boxed{
  \dot\rho = -\ii[H,\rho]
  + \sum_k \Bigl(L_k\rho L_k^\dagger
    - \tfrac{1}{2}\bigl\{L_k^\dagger L_k,\rho\bigr\}\Bigr).
  }
  \label{eq:Lindblad}
\end{equation}
The \textbf{jump operators} $L_k$ encode individual decay channels.

\subsection{Physical decay channels}

\begin{table}[h!]
\centering
\renewcommand{\arraystretch}{1.4}
\begin{tabular}{llll}
\toprule
Jump operator & Rate & Process & Timescale \\
\midrule
$L_1 = \sqrt{\gamma_q}\,\smq$ & $\gamma_q = 1/T_1^{(q)}$ & Qubit energy relaxation & $T_1^{(q)}$ \\
$L_2 = \sqrt{\gamma_t}\,\smt$ & $\gamma_t = 1/T_1^{(t)}$ & TLS energy relaxation & $T_1^{(t)}$ \\
$L_3 = \sqrt{\gamma_{\phi,q}/2}\,\szq$ & $\gamma_{\phi,q} = 2/T_\phi^{(q)}$ & Qubit pure dephasing & $T_\phi^{(q)}$ \\
$L_4 = \sqrt{\gamma_{\phi,t}/2}\,\szt$ & $\gamma_{\phi,t} = 2/T_\phi^{(t)}$ & TLS pure dephasing & $T_\phi^{(t)}$ \\
\bottomrule
\end{tabular}
\caption{Jump operators for the qubit--TLS system. The total qubit
dephasing rate is $1/T_2^{(q)} = 1/(2T_1^{(q)}) + 1/T_\phi^{(q)}$,
and similarly for the TLS.}
\label{tab:jump}
\end{table}

\subsection{Equations of motion for the density matrix elements}

Writing $\rho$ in the basis $\{\ket{ee},\ket{eg},\ket{ge},\ket{gg}\}$,
\eqref{eq:Lindblad} yields $16$ coupled ODEs (reduced to $12$ independent
equations by Hermiticity and $\Tr\rho = 1$). In the single-excitation
subspace, the relevant elements are
$\rho_{eg,eg}$, $\rho_{ge,ge}$, and $\rho_{eg,ge}$ (coherence).
Their equations of motion are
\begin{align}
  \dot\rho_{eg,eg}
    &= \ii g(\rho_{ge,eg} - \rho_{eg,ge})
      - \gamma_q\,\rho_{eg,eg}
      + \gamma_t\,\rho_{ee,ee},
      \label{eq:rho_egeg}\\
  \dot\rho_{ge,ge}
    &= \ii g(\rho_{eg,ge} - \rho_{ge,eg})
      - \gamma_t\,\rho_{ge,ge}
      + \gamma_q\,\rho_{eg,eg},
      \label{eq:rho_gege}\\
  \dot\rho_{eg,ge}
    &= \ii g(\rho_{ge,ge} - \rho_{eg,eg})
      - \ii\D\,\rho_{eg,ge}
      - \frac{\gamma_q + \gamma_t
              + \gamma_{\phi,q} + \gamma_{\phi,t}}{2}\,\rho_{eg,ge}.
      \label{eq:rho_egge}
\end{align}
These equations are solved numerically using QuTiP \cite{Johansson2012}.

% ════════════════════════════════════════════════════════════════════════════
\section{The Rotating Frame}

Working in the lab frame, all quantities oscillate rapidly at $\wq$.
It is conventional to move to a \textbf{rotating frame} via the unitary
\begin{equation}
  R(t) = \exp\!\Bigl(\ii\frac{\wd}{2}(\szq + \szt)\,t\Bigr),
\end{equation}
where $\wd$ is a reference (often drive) frequency.
The transformed Hamiltonian $\tilde H = R^\dagger H R - \ii R^\dagger\dot R$
in the rotating frame (with $\wd = \wq$ and $\D = \wq - \wt$) is
\begin{equation}
  \tilde H = -\frac{\D}{2}\szt
           + g\bigl(\spq\smt + \smq\spt\bigr).
  \label{eq:H_rot}
\end{equation}
This frame removes the fast $\wq$ oscillation and is the natural
frame for numerical simulations and for applying the dispersive
approximation \cite{Blais2021}.

% ════════════════════════════════════════════════════════════════════════════
\section{The Dispersive Limit and Schrieffer--Wolff Transformation}

When $|\D| \gg g$, direct excitation exchange is suppressed.
A \textbf{Schrieffer--Wolff} (SW) unitary transformation
$U_{\mathrm{SW}} = \exp(\lambda S)$ with
$\lambda = g/\D$ and $S = \spq\smt - \smq\spt$
perturbatively eliminates the off-diagonal coupling to second order in
$g/\D$ \cite{Blais2021,Bravyi2011}. The effective Hamiltonian is
\begin{equation}
  H_{\mathrm{eff}}
  = \frac{\wq}{2}\szq + \frac{\wt}{2}\szt
  + \chi\,\szq\szt,
  \qquad
  \chi = \frac{g^2}{\D}.
  \label{eq:H_disp}
\end{equation}
The $\chi\,\szq\szt$ term is the \textbf{dispersive interaction}: the
qubit and TLS shift each other's frequency depending on their mutual
state, without exchanging excitations. This is the operating regime for
qubit readout in circuit QED \cite{Blais2021}.

% ════════════════════════════════════════════════════════════════════════════
\section{Connection to the Solomon Equations}

The Solomon equations \cite{Solomon1955,Ashhab2006} describe the
qubit--TLS dynamics in the \emph{classical rate-equation limit}, where
all quantum coherences are discarded. They read
\begin{align}
  \dot P_q &= -\gamma_q P_q
             - \Gamma(P_q - P_t), \label{eq:Solomon_q}\\
  \dot P_t &= -\gamma_t P_t
             + \Gamma(P_q - P_t), \label{eq:Solomon_t}
\end{align}
where $P_q = \langle\szq\rangle/2 + 1/2$ is the qubit excited-state
population, $P_t$ is the TLS excited-state population, and
$\Gamma = 2g^2\gamma_t/(\D^2 + \gamma_t^2)$ is the
\textbf{exchange rate} (Lorentzian in the detuning $\D$).

The Solomon equations emerge from the full Lindblad equation
\eqref{eq:Lindblad} by:
\begin{enumerate}
  \item Moving to the secular approximation (averaging over fast Rabi
        oscillations at $\Omega_R$).
  \item Setting all off-diagonal density matrix elements (coherences) to
        zero: $\rho_{eg,ge} = 0$.
  \item This is valid when $\Omega_R \gg \gamma_q, \gamma_t$
        (strong coupling) or $|\D| \gg g$ (dispersive, where coherences
        decay faster than populations).
\end{enumerate}
The key comparison to make in this thesis is therefore:

\begin{center}
\renewcommand{\arraystretch}{1.5}
\begin{tabular}{lll}
\toprule
Model & Equations & Validity \\
\midrule
Full quantum (Lindblad) & Eqs.~\eqref{eq:rho_egeg}--\eqref{eq:rho_egge} & Always (within RWA) \\
Solomon (rate equations) & Eqs.~\eqref{eq:Solomon_q}--\eqref{eq:Solomon_t}
  & $\Omega_R \gg \gamma_q,\gamma_t$ or $|\D| \gg g$ \\
\bottomrule
\end{tabular}
\end{center}

The Solomon equations are expected to break down when
$g \sim \gamma_q$ or $g \sim \gamma_t$ (strong coupling to bath), or
when $T_2^{(\mathrm{TLS})}$ is long enough that coherences survive.
Mapping out this breakdown numerically is \textbf{Task 3} of the thesis.

% ════════════════════════════════════════════════════════════════════════════
\section{Summary of the Derivation Roadmap}

\begin{enumerate}
  \item \textbf{Bare Hamiltonians} $H_q$, $H_{\mathrm{TLS}}$
        \eqref{eq:Hq}--\eqref{eq:Ht}: two independent spins.
  \item \textbf{Full coupling} $H_{\mathrm{int}}^{\mathrm{full}}$
        \eqref{eq:Hint_full}: all four bilinear terms.
  \item \textbf{RWA} \eqref{eq:Hint_RWA}: drop counter-rotating terms
        when $g \ll \wq$.
  \item \textbf{Total Hamiltonian} \eqref{eq:H_total} with block-diagonal
        structure \eqref{eq:H_block}.
  \item \textbf{Diagonalization} \eqref{eq:eigenvalues}--\eqref{eq:psi_t}:
        dressed states, Rabi oscillations (closed system).
  \item \textbf{Open system}: promote to density matrix, add jump operators
        (Table~\ref{tab:jump}), write Lindblad equation \eqref{eq:Lindblad}.
  \item \textbf{Rotating frame} \eqref{eq:H_rot}: remove fast oscillations.
  \item \textbf{Dispersive limit} \eqref{eq:H_disp}: Schrieffer--Wolff
        transformation for $|\D| \gg g$.
  \item \textbf{Solomon limit} \eqref{eq:Solomon_q}--\eqref{eq:Solomon_t}:
        discard coherences, identify breakdown conditions.
\end{enumerate}

% ════════════════════════════════════════════════════════════════════════════
\section*{Acknowledgements}
The author thanks Mathieu Fechant and Nico Gosling for guidance on this
project, and Ioan Pop for arranging the thesis placement at KIT.

% ════════════════════════════════════════════════════════════════════════════
\bibliographystyle{unsrtnat}
\begin{thebibliography}{99}

\bibitem{Blais2021}
A.~Blais, A.~L. Grimsmo, S.~M. Girvin, and A.~Wallraff,
\emph{Circuit quantum electrodynamics},
Rev.\ Mod.\ Phys.\ \textbf{93}, 025005 (2021).

\bibitem{Haroche2006}
S.~Haroche and J.-M.~Raimond,
\emph{Exploring the Quantum: Atoms, Cavities, and Photons},
Oxford University Press (2006).

\bibitem{WallsMilburn2008}
D.~F. Walls and G.~J. Milburn,
\emph{Quantum Optics}, 2nd ed.,
Springer (2008).

\bibitem{Breuer2002}
H.-P. Breuer and F.~Petruccione,
\emph{The Theory of Open Quantum Systems},
Oxford University Press (2002).

\bibitem{Gardiner2004}
C.~W. Gardiner and P.~Zoller,
\emph{Quantum Noise}, 3rd ed.,
Springer (2004).

\bibitem{Ashhab2006}
S.~Ashhab, J.~R. Johansson, and F.~Nori,
\emph{Decoherence in a superconducting quantum bit circuit},
Phys.\ Rev.\ B \textbf{74}, 184415 (2006).
[arXiv:cond-mat/0604475]

\bibitem{Muller2019}
C.~M\"uller, J.~H. Cole, and J.~Lisenfeld,
\emph{Towards understanding two-level-systems in amorphous solids},
Rep.\ Prog.\ Phys.\ \textbf{82}, 124501 (2019).

\bibitem{Lisenfeld2019}
J.~Lisenfeld \emph{et al.},
\emph{Electric field spectroscopy of material defects in transmon qubits},
npj Quantum Inf.\ \textbf{5}, 105 (2019).

\bibitem{Solomon1955}
I.~Solomon,
\emph{Relaxation processes in a system of two spins},
Phys.\ Rev.\ \textbf{99}, 559 (1955).

\bibitem{Bravyi2011}
S.~Bravyi, D.~P. DiVincenzo, and D.~Loss,
\emph{The Schrieffer--Wolff transformation for quantum many-body systems},
Ann.\ Phys.\ \textbf{326}, 2793 (2011).

\bibitem{Johansson2012}
J.~R. Johansson, P.~D. Nation, and F.~Nori,
\emph{QuTiP: An open-source Python framework for the dynamics of open
quantum systems},
Comp.\ Phys.\ Comm.\ \textbf{183}, 1760 (2012).

\bibitem{Kam2025}
J.~F. Kam \emph{et al.},
\emph{Detrimental non-Markovian errors for surface code memory},
arXiv:2507.13779 (2025).

\end{thebibliography}

\end{document}
